\documentclass[11pt]{beamer}
%\usepackage{CJKutf8}
\usepackage{beamerthemesplit}
\usepackage{}
%\usetheme{Frankfurt}
%\usetheme{CambridgeUS}
%\usecolortheme{beaver}
%\usetheme{AnnArbor}
\usetheme{Dresden}
%\usebeamercolor{beetle}
\usepackage{xeCJK}
\setCJKmainfont{AR PL KaitiM GB}
%\useoutertheme{miniframes}
\usepackage{amsmath}
\usepackage{graphicx}
\usepackage{float} 
\usepackage{subfigure}
\usepackage{amssymb}
\usepackage{graphicx}
\usepackage{eufrak}
\usepackage{color}
\usepackage{array}
\usepackage{slashed}
\usepackage{simplewick}
\usepackage{tikz}
\usepackage{tcolorbox}
\usepackage[T1]{fontenc}
\graphicspath{{../figures/}}

%%figures
\def\lfig#1#2{\includegraphics[width=#1 in]{#2}}
\def\addfig#1#2{\begin{center}\includegraphics[width=#1 in]{#2}\end{center}}
\def\wulian{\includegraphics[width=0.18in]{emoji_wulian.jpg}}
\def\bigwulian{\includegraphics[width=0.35in]{emoji_wulian.jpg}}
\def\bye{\includegraphics[width=0.18in]{emoji_bye.jpg}}
\def\bigbye{\includegraphics[width=0.35in]{emoji_bye.jpg}}
\def\huaixiao{\includegraphics[width=0.18in]{emoji_huaixiao.jpg}}
\def\bighuaixiao{\includegraphics[width=0.35in]{emoji_huaixiao.jpg}}
\def\jianxiao{\includegraphics[width=0.18in]{emoji_jianxiao.jpg}}
\def\bigjianxiao{\includegraphics[width=0.35in]{emoji_jianxiao.jpg}}
%% colors
\def\blacktext#1{{\color{black}#1}}
\def\bluetext#1{{\color{blue}#1}}
\def\redtext#1{{\color{red}#1}}
\def\darkbluetext#1{{\color[rgb]{0,0.2,0.6}#1}}
\def\skybluetext#1{{\color[rgb]{0.2,0.7,1.}#1}}
\def\cyantext#1{{\color[rgb]{0.,0.5,0.5}#1}}
\def\greentext#1{{\color[rgb]{0,0.7,0.1}#1}}
\def\darkgray{\color[rgb]{0.2,0.2,0.2}}
\def\lightgray{\color[rgb]{0.6,0.6,0.6}}
\def\gray{\color[rgb]{0.4,0.4,0.4}}
\def\blue{\color{blue}}
\def\red{\color{red}}
\def\green{\color{green}}
\def\darkgreen{\color[rgb]{0,0.4,0.1}}
\def\darkblue{\color[rgb]{0,0.2,0.6}}
\def\skyblue{\color[rgb]{0.2,0.7,1.}}
%%control
\def\be{\begin{equation}}
\def\ee{\nonumber\end{equation}}
\def\bea{\begin{eqnarray*}}
\def\eea{\nonumber\end{eqnarray*}}
\def\bch{}
\def\ech{}
\def\bitem{\begin{itemize}}
\def\eitem{\end{itemize}}
\def\bcenter{\begin{center}}
\def\ecenter{\end{center}}
\def\bex{\begin{minipage}{0.2\textwidth}\includegraphics[width=0.6in]{jugelizi.png}\end{minipage}\begin{minipage}{0.76\textwidth}}
\def\eex{\end{minipage}}
\def\chtitle#1{\frametitle{\bch#1\ech}}
\def\bmat#1{\left(\begin{array}{#1}}
\def\emat{\end{array}\right)}
\def\bcase#1{\left\{\begin{array}{#1}}
\def\ecase{\end{array}\right.}
\def\bmini#1{\begin{minipage}{#1\textwidth}}
\def\emini{\end{minipage}}
\def\tbox#1{\begin{tcolorbox}#1\end{tcolorbox}}
\def\pfrac#1#2#3{\left(\frac{\partial #1}{\partial #2}\right)_{#3}}
%%symbols
\def\bropt{\,(\ \ \ )}
\def\sone{$\star$}
\def\stwo{$\star\star$}
\def\sthree{$\star\star\star$}
\def\sfour{$\star\star\star\star$}
\def\sfive{$\star\star\star\star\star$}
\def\rint{{\int_\leftrightarrow}}
\def\roint{{\oint_\leftrightarrow}}
\def\stdHf{{\textit{\r H}_f}}
\def\deltaH{{\Delta \textit{\r H}}}
\def\ii{{\dot{\imath}}}
\def\skipline{{\vskip0.1in}}
\def\skiplines{{\vskip0.2in}}
\def\lagr{{\mathcal{L}}}
\def\hamil{{\mathcal{H}}}
\def\vecv{{\mathbf{v}}}
\def\vecx{{\mathbf{x}}}
\def\vecy{{\mathbf{y}}}
\def\veck{{\mathbf{k}}}
\def\vecp{{\mathbf{p}}}
\def\vecn{{\mathbf{n}}}
\def\vecA{{\mathbf{A}}}
\def\vecP{{\mathbf{P}}}
\def\vecsigma{{\mathbf{\sigma}}}
\def\hatJn{{\hat{J_\vecn}}}
\def\hatJx{{\hat{J_x}}}
\def\hatJy{{\hat{J_y}}}
\def\hatJz{{\hat{J_z}}}
\def\hatj#1{\hat{J_{#1}}}
\def\hatphi{{\hat{\phi}}}
\def\hatq{{\hat{q}}}
\def\hatpi{{\hat{\pi}}}
\def\vel{\upsilon}
\def\Dint{{\mathcal{D}}}
\def\adag{{\hat{a}^\dagger}}
\def\bdag{{\hat{b}^\dagger}}
\def\cdag{{\hat{c}^\dagger}}
\def\ddag{{\hat{d}^\dagger}}
\def\hata{{\hat{a}}}
\def\hatb{{\hat{b}}}
\def\hatc{{\hat{c}}}
\def\hatd{{\hat{d}}}
\def\hatN{{\hat{N}}}
\def\hatH{{\hat{H}}}
\def\hatp{{\hat{p}}}
\def\Fup{{F^{\mu\nu}}}
\def\Fdown{{F_{\mu\nu}}}
\def\newl{\nonumber \\}
\def\vece{\mathrm{e}}
\def\calM{{\mathcal{M}}}
\def\calT{{\mathcal{T}}}
\def\calR{{\mathcal{R}}}
\def\barpsi{\bar{\psi}}
\def\baru{\bar{u}}
\def\barv{\bar{\upsilon}}
\def\qeq{\stackrel{?}{=}}
\def\torder#1{\mathcal{T}\left(#1\right)}
\def\rorder#1{\mathcal{R}\left(#1\right)}
\def\contr#1#2{\contraction{}{#1}{}{#2}#1#2}
\def\trof#1{\mathrm{Tr}\left(#1\right)}
\def\trace{\mathrm{Tr}}
\def\comm#1{\ \ \ \left(\mathrm{used}\ #1\right)}
\def\tcomm#1{\ \ \ (\text{#1})}
\def\slp{\slashed{p}}
\def\slk{\slashed{k}}
\def\calp{{\mathfrak{p}}}
\def\veccalp{\mathbf{\mathfrak{p}}}
\def\Tthree{T_{\tiny \textcircled{3}}}
\def\pthree{p_{\tiny \textcircled{3}}}
\def\dbar{{\,\mathchar'26\mkern-12mu d}}
\def\erf{\mathrm{erf}}
\def\const{\mathrm{constant}}
\def\pheat{\pfrac p{\ln T}V}
\def\vheat{\pfrac V{\ln T}p}
%%units
\def\fdeg{{^\circ \mathrm{F}}}
\def\cdeg{^\circ \mathrm{C}}
\def\atm{\,\mathrm{atm}}
\def\angstrom{\,\text{\AA}}
\def\SIL{\,\mathrm{L}}
\def\SIkm{\,\mathrm{km}}
\def\SIyr{\,\mathrm{yr}}
\def\SIGyr{\,\mathrm{Gyr}}
\def\SIV{\,\mathrm{V}}
\def\SImV{\,\mathrm{mV}}
\def\SIeV{\,\mathrm{eV}}
\def\SIkeV{\,\mathrm{keV}}
\def\SIMeV{\,\mathrm{MeV}}
\def\SIGeV{\,\mathrm{GeV}}
\def\SIcal{\,\mathrm{cal}}
\def\SIkcal{\,\mathrm{kcal}}
\def\SImol{\,\mathrm{mol}}
\def\SIN{\,\mathrm{N}}
\def\SIHz{\,\mathrm{Hz}}
\def\SIm{\,\mathrm{m}}
\def\SIcm{\,\mathrm{cm}}
\def\SIfm{\,\mathrm{fm}}
\def\SImm{\,\mathrm{mm}}
\def\SInm{\,\mathrm{nm}}
\def\SImum{\,\mathrm{\mu m}}
\def\SIJ{\,\mathrm{J}}
\def\SIW{\,\mathrm{W}}
\def\SIkJ{\,\mathrm{kJ}}
\def\SIs{\,\mathrm{s}}
\def\SIkg{\,\mathrm{kg}}
\def\SIg{\,\mathrm{g}}
\def\SIK{\,\mathrm{K}}
\def\SImmHg{\,\mathrm{mmHg}}
\def\SIPa{\,\mathrm{Pa}}
\def\secpage#1#2{\begin{frame}\bch\bcenter{\bf \Huge #1} \skipline \tbox{\bcenter #2\ecenter}\ecenter\ech\end{frame}}

\title{Cosmology}
  \author{Haoting}
  \date{\today}
\begin{document}

\begin{frame}
\begin{center}
\begin{Large}
 {\bf I}ntroductory  {\bf C}osmology 
{\vskip 0.1in}
Lecture 02-The Dynamics of Spacetime
\end{Large}
\end{center}
\vskip 0.1in
\begin{center}
Haoting Xu
\vskip 0.1in
xuht9@mail2.sysu.edu.cn
\vskip 0.1in
{\tiny \url{https://github.com/HaotingXu/seminar_lec/} }\\
\end{center}
\end{frame}

\section{Intro}
\begin{frame}\frametitle{宇宙如何膨胀}
上一讲我们提到，在宇宙学中，时空几何使用FRW度规来描述。
\be 
ds^2 = -c^2dt^2+ a^2(t) \lmbk{\frac{1}{1-kr^2/R^2}dr^2+r^2\lbk{d\theta^2+\sin^2\theta d\phi^2}}
\ee 
这只是一个对均匀各项同性的宇宙学原理的数学表述，并没有任何物理。物理是来决定这个度规的动力学，即$a(t)$如何演化。这一讲我们就来求解$a(t)$如何随时间演化。
\end{frame}
\begin{frame}\frametitle{广义相对论}
决定$a(t)$如何演化的物理当然是广义相对论。广义相对论可以简单理解为

Spacetime tells matter how to move, matter tells space how to curve. 

于是，几何的部分我们已经有了，而另一部分我们要解决的就是物质，或者更精确地说，物质的能量动量张量。
\end{frame}
\section{Perfect Fluids}
\secpage{Perfect Fluids}{$\rho(t) = \rho_0 a^{-3(1+w)}$}
\begin{frame}\frametitle{理想流体}
对于我们人类来说，一个星系的体积十分巨大。但是如果我们从大尺度来思考问题，每个星系实际上相当于宇宙流体的一个小原子。

对于理想流体我们只关心两件事情：第一个是流体的能量密度$\rho(t)$(在相对论的文本中，质量和能量不分家)，第二个是流体的压强$P(t)$。
\end{frame}
\begin{frame}\frametitle{研究对象}
在平衡态的情况下，$P(\rho)$被称作状态方程。我们主要关心三种东西的状态方程。
\begin{itemize}
	\item 非相对论气体：尘埃。
	\item 相对论气体。
	\item 暗能量。
\end{itemize}
\end{frame}
\begin{frame}\frametitle{非相对论气体}
在这里我们直接采用统计物理中的结果，如果不熟悉的话可以去看我的统计物理讲义。对于非相对论(理想)气体，压强为
\be 
P = \frac{Nk_BT}{V}  = \frac{1}{3} \frac{N}{V} m\meanvl{v^2}
\ee 
这里用到了理想气体的速度平方的平均值。因为是非相对论情况，动能远远小于静止的能量$mc^2$，代入$E\simeq mc^2$，得到
\be 
P = \frac{NE}{V}\frac{\meanvl{v^2}}{c^2} \simeq 0
\ee 
所以我们平时讲的那点压强，在宇宙学尺度看来，几乎为$0$。
\end{frame}
\begin{frame}\frametitle{相对论气体}
对于极端相对论情况，我们在统计物理中推导过，压强和能量密度的关系为
\be 
P = \frac{1}{3} \rho 
\ee 
大部分的状态方程都有如下形式
\be 
P = w\rho 
\ee 
我们已经看到非相对论气体$w=0$，相对论气体$w = 1/3$。我们以后将会看到一些不可思议的例子。
\end{frame}
\begin{frame}\frametitle{对$w$的一个理论上的限制}
可以从流体力学中推导出来，介质中的声速满足如下关系
\be 
c_s^2 =c^2 \frac{dP}{d\rho}
\ee 
这告诉我们$w \le 1$，因为声速不能超过光速。$w$如果是负的，这时候声速$c_s$变成一个虚数，这意味着在这种介质中，声波不可能传播，因为它会指数衰减。
\end{frame}
\begin{frame}{连续性方程(4-动量流守恒方程)}
下面介绍一个决定物质的能量密度如何演化的方程。这个方程的严谨版本需要从广义相对论出发。我们之后会介绍一下这个思想。下面我们给出非严谨版本的推导，但是这个经典的随便的推导竟然和严谨的一样。考虑没有热量的热力学第一定律
\be 
dE = -p dV
\ee 
则有
\be 
\frac{dE}{dt} = -p \frac{dV}{dt}
\ee
假设今天这团东西的体积是$V_0$，那么在任何时间有$V = a^3(t)V_0$，$E = \rho a^3(t) V_0$。直接将这两个式子代入上式，得到
\be 
\dot{\rho} a^3 + 3\rho a^2 \dot{a} = -3pa^2 \dot{a}
\ee 
利用哈勃参数的定义，有
\be 
\dot{\rho} + 3H(\rho+P) =0 
\ee 
\end{frame}
\begin{frame}\frametitle{求解物质的演化}
把一般的状态方程$P = w\rho$代入上式，得到
\be 
\frac{\dot{\rho}}{\rho} = -3(1+w) \frac{\dot{a}}{a}
\ee 
直接求解，得到
\be 
\rho (t) = \rho_0 a^{-3(1+w)}
\ee 
其中$\rho_0$是今天的能量密度。值得记住两个例子，对于尘埃，有
\be 
\rho_m \sim \frac{1}{a^3}
\ee 
这和我们瞎猜的差不多一样。再看辐射
\be 
\rho_r \sim \frac{1}{a^4}
\ee 
这也和我们瞎猜的差不多一样，相比于尘埃，多了一个$1/a$因子，这恰好就是红移。
\end{frame}
\section{Friedmann Equation}
\secpage{Friedmann Equation}{$\lbk{\frac{\dot{a}}{a}}^2 = \frac{8\pi G}{3c^2}\rho - \frac{kc^2}{R^2a^2}$}
\begin{frame}\frametitle{Friedmann 方程}
重复一下我们的图像：我们考虑的是一个均匀(里面的物质是均匀的)且各向同性的宇宙。在FRW度规下和上面所述的物质的背景下，爱因斯坦方程就变成了
\be 
H^2 \equiv \fbk{\dot{a}}{a}^2 = \frac{8\pi G}{3c^2} \rho - \frac{kc^2}{R^2a}^2
\ee 
其中$k$是空间曲率，$G = 6.67\times 10^{-11} \mathrm{m^3 Kg^{-1} s^{-2}}$。因为我们没学过广义相对论，所以这部分推导暂时没法推。但是，出乎意料的是，如果做一些错误的(不完全符合)图像的假设，用牛顿引力也能推导出上面的方程。但是读者应该牢记，这个推导是“假”的。有关正确的推导，大家可以去看Dodelson的书。
\end{frame}
\begin{frame}\frametitle{求解宇宙}
最后，有了Friedmann方程和物质的演化方程(动力学方程)和物质的状态方程，就可以求解宇宙是如何演化的了。

这里就和静电场非常相似，要求解一个具体问题，首先要知道一个问题的源和电场遵从的方程(泊松方程)。
\end{frame}
\begin{frame}\frametitle{引力场}
物理学家很喜欢“场”的概念。干什么都要先定义一波场。类比于静电场，我们定义引力场，用一个引力势来描述
\be 
\vec{F} = -m\nabla \Phi  = - \int_{V^{\prime}} (\rho dV^{\prime}) \frac{Gm}{|\vec{x}-\vec{x}^{\prime}|^2}
\ee 
其中上式最后一个等号用到了牛顿万有引力定律。直接对上式求散度，则有
\bea 
-\nabla \cdot (\nabla \Phi) &=& -\nabla \cdot \lbk{\int (\rho dV^{\prime}) \frac{G}{|\vec{x}-\vec{x}^{\prime}|^2}}\\
&=& -\int (\rho dV^{\prime}) G 4\pi \delta^{(3)} (vec{x}-\vec{x}^{\prime}) \\
&=& -4\pi G\rho (\vec{x})
\eea 
这就是引力场的泊松方程。最后，我们把$\rho$换成能量密度(原来是质量密度)，则有
\be 
\nabla^2 \Phi = \frac{4\pi G}{c^2} \rho 
\ee 
\end{frame}
\begin{frame}\frametitle{求解场}
我们考虑如下模型，整个宇宙就是一个大球，取$m$为大球边缘的一个小质量元，设整个大球的质量为$M(r) = \rho 4\pi r^3/3$。我们假设它是不变的。对泊松方程积分一次得到高斯定理
\be 
\int_V \nabla \cdot (\nabla \Phi) dV = \int_S (\nabla \Phi)\cdot d\vec{S} = (\nabla \Phi)_r \cdot 4\pi r^2= \int_V \frac{4\pi G}{c^2} \rho dV
\ee 
求解得到
\be 
(\nabla \Phi)_r = \frac{GM(r)}{r^2}
\ee 
于是我们有
\be 
m\ddot{r} = -\frac{GMm}{r^2}
\ee 
\end{frame}
\begin{frame}\frametitle{对动力学方程积分}
把上式对时间积分一次，得到
\be 
\int m\ddot{r}\dot{r}dt = - \int \frac{GMm}{r^2} \dot{r}dt 
\ee 
第一项分布积分，得到
\be 
\frac{1}{2} m\dot{r}^2 - \frac{GMm}{r} =E
\ee
上式就是大名鼎鼎的能量守恒定律。我认为上面这一套是物理系学生的看家本领。
\end{frame}
\begin{frame}\frametitle{Friedmann方程}
根据之前的FRW度规，代入$r= a(t)r_0$，我们定义
\be 
C = -\frac{2E}{r_0^2}
\ee 
得到Friedmann方程
\be 
\fbk{\dot{a}}{a} = \frac{8\pi G}{c^2} -\frac{C}{a^2}
\ee 
\end{frame}
\begin{frame}\frametitle{曲率}
在这里我们强行说$ C =  k^2c^2/R^2$。这是牛顿力学无法理解的。但是可以从牛顿力学有一个很好的类比。
\begin{itemize}
\item 当$C<0$时，这意味着$E>0$，而我们知道在两体运动当中，$E>0$情况对应于双曲线。双曲线就意味着$k$为负，是负曲率空间。我们将在以后看到，负曲率的宇宙的$a(t)$是大概率发散的。
\item 当$C>0$时，这意味着$E<0$，而在牛顿力学，$E<0$对应着一条闭合曲线。而我们以后会看到，正曲率的宇宙大概率是收敛的。
\end{itemize}
\end{frame}
\begin{frame}\frametitle{错误的图像给出了正确的结果}
注意我们上面推导中的图像是错误的，宇宙不是一个巨大的物质球。而且这样宇宙膨胀的图像也是错误的，宇宙并不是以某一点为中心而膨胀，而是宇宙中有无限的空间，空间中有一个个网格，是那个网格(以$a$来表示尺度)在膨胀。一个星系的共动坐标变化并不是“宇宙”意义上的膨胀。事实上从哈勃图我们也能发现，星系的共动坐标变化相对于宇宙膨胀的速度是几乎可以忽略不计的。
\end{frame}


\end{document}
